\documentclass[a4paper,12pt]{ctexart}
\usepackage{xcolor}
\usepackage{graphicx}
\usepackage{amsmath}
\usepackage{hyperref}
\hypersetup{colorlinks=true,linkcolor=blue!70!black}
\usepackage[top=1.8cm, bottom=1.8cm, left=1.8cm, right=1.8cm]{geometry}
\usepackage{array}
\linespread{1.35}

\definecolor{earthgreen}{RGB}{30,120,70}
\definecolor{skyblue}{RGB}{0,100,180}
\definecolor{thinkorange}{RGB}{230,120,20}

\newenvironment{thinkbox}{
  \vspace{0.25cm}
  \noindent
  \begin{minipage}{0.92\textwidth}
  \colorbox{blue!4}{
  \begin{minipage}{0.94\textwidth}
  \vspace{0.1cm}
  \textbf{\textcolor{thinkorange}{【想一想】}}
}{
  \vspace{0.1cm}
  \end{minipage}}
  \end{minipage}
  \par
  \vspace{0.25cm}
}

\newenvironment{knowbox}[1]{
  \vspace{0.2cm}
  \noindent\colorbox{green!5}{\parbox{\dimexpr\linewidth-2\fboxsep}{\textbf{\textcolor{earthgreen}{知识点：}} #1}}\par
  \vspace{0.2cm}
}

\title{\textcolor{skyblue}{\Huge\textbf{地理初步}}}
\author{地球与宇宙 $\cdot$ 地形与气候 $\cdot$ 中国与世界}
\date{}

\begin{document}
\maketitle
\thispagestyle{empty}

\section*{\textcolor{skyblue}{给同学的话}}

地理是一门"上知天文、下知地理"的学科。它研究的是我们脚下的地球——它的形状、它的运动、它表面的山川河流、它头顶的大气层。

你可能会问：这些知识有什么用？
\begin{itemize}
  \item 知道地球的形状和运动，你才能理解为什么有时差、为什么有四季
  \item 知道气候的分布，你才能理解为什么有的地方是沙漠、有的地方是雨林
  \item 知道地形特点，你才能理解为什么地震和火山集中在某些地区
\end{itemize}

\vspace{0.3cm}

\textbf{本模块分为两大部分：}
\begin{itemize}
  \item \textbf{自然地理}（分册一：自然地理基础）—— 地球、地图、地形、气候
  \item \textbf{灾害与资源地理}（分册二：地貌、灾害与资源环境）—— 板块构造、内外力作用、自然灾害、自然资源
\end{itemize}

\newpage

\section{\textcolor{orange!80!black}{第一课\quad 地球与宇宙}}

\subsection*{1.1\quad 地球的形状和大小}

古人以为天圆地方，后来人们发现地球其实是一个\textbf{球体}。最早证明这一点的是古希腊人——他们看到远航的船只先消失船身、再消失桅杆，说明海面是弯曲的。

\begin{knowbox}{地球的基本数据}
  \begin{itemize}
    \item 赤道半径：约 6378 千米
    \item 极半径：约 6357 千米（地球是略扁的椭球体）
    \item 赤道周长：约 40076 千米
    \item 表面积：约 5.1 亿平方公里
  \end{itemize}
\end{knowbox}

\begin{thinkbox}
地球为什么是"两头扁、中间鼓"的形状？回想一下我们学过的离心力——地球自转把赤道部分"甩"出去了。
\end{thinkbox}

\subsection*{1.2\quad 地球的自转与公转}

\textbf{自转}：地球绕自身轴旋转。
\begin{itemize}
  \item 方向：自西向东
  \item 周期：约 24 小时（一天）
  \item 产生：昼夜交替、太阳东升西落
\end{itemize}

\textbf{公转}：地球绕太阳运行。
\begin{itemize}
  \item 方向：自西向东
  \item 周期：约 365.25 天（一年）
  \item 地轴倾斜：约 $23.5^\circ$——这是产生\textbf{四季}的原因
\end{itemize}

\subsection*{1.3\quad 经纬度}

\begin{knowbox}{
  \textbf{经线（子午线）}：连接南北两极的半圆，指示南北方向。\\
  \textbf{纬线}：与赤道平行的圆圈，指示东西方向。\\
  \textbf{赤道}：$0^\circ$纬线，把地球分为南北半球。\\
  \textbf{本初子午线}：$0^\circ$经线，通过伦敦格林尼治天文台。}
\end{knowbox}

\begin{thinkbox}
北京的位置大约是北纬 $40^\circ$、东经 $116^\circ$。你能在地球仪上找到这个位置吗？和你所在的城市相比，谁的纬度更高？
\end{thinkbox}

\newpage

\section{\textcolor{orange!80!black}{第二课\quad 地图阅读}}

地图不是现实的"照片"，而是现实的"简化模型"。地图有三大要素：

\subsection*{2.1\quad 比例尺}

比例尺 = 图上距离 $\div$ 实际距离。
\begin{itemize}
  \item 大比例尺（如 $1:10000$）：表示范围小，内容详细
  \item 小比例尺（如 $1:10000000$）：表示范围大，内容简略
\end{itemize}

\subsection*{2.2\quad 方向}

\begin{itemize}
  \item 一般地图：上北下南、左西右东
  \item 有指向标的地图：指向标箭头指向北
  \item 有经纬网的地图：经线指示南北、纬线指示东西
\end{itemize}

\subsection*{2.3\quad 图例和注记}

\begin{center}
\begin{tabular}{|c|c||c|c|}
\hline
\textbf{图例} & \textbf{含义} & \textbf{图例} & \textbf{含义} \\ \hline
$\bullet$ & 首都 & $\sim\sim\sim$ & 河流 \\
$\bigcirc$ & 一般城市 & $\cdots\cdots$ & 国界 \\
$\blacktriangle$ & 山峰 & $=$ & 铁路 \\
\hline
\end{tabular}
\end{center}

\begin{thinkbox}
现在你手机上随时都能打开电子地图，但电子地图和纸质地图有一个根本区别：纸质地图是一张"固定"的图，而电子地图可以根据需要放大缩小。你知道电子地图的"缩放"背后用的是什么技术吗？
\end{thinkbox}

\newpage

\section{\textcolor{orange!80!black}{第三课\quad 地形与地貌}}

地球表面高低起伏，主要有五种基本地形：

\subsection*{五种地形类型}

\begin{center}
\begin{tabular}{|c|c|c|c|}
\hline
\textbf{地形} & \textbf{海拔} & \textbf{起伏} & \textbf{例子} \\ \hline
山地 & $>500$米 & 起伏很大 & 喜马拉雅山 \\
高原 & $>500$米 & 起伏不大 & 青藏高原 \\
平原 & $<200$米 & 平坦开阔 & 华北平原 \\
丘陵 & $<500$米 & 起伏和缓 & 东南丘陵 \\
盆地 & 四周高中间低 & — & 四川盆地 \\
\hline
\end{tabular}
\end{center}

\subsection*{中国地形概览}

中国的地形特征是"\textbf{西高东低，呈三级阶梯分布}"：
\begin{itemize}
  \item \textbf{第一级阶梯}：青藏高原（平均海拔 4000 米以上）
  \item \textbf{第二级阶梯}：内蒙古高原、黄土高原、云贵高原、四川盆地（1000--2000 米）
  \item \textbf{第三级阶梯}：东北平原、华北平原、长江中下游平原（500 米以下）
\end{itemize}

\begin{thinkbox}
中国的地形"西高东低"对河流和气候有什么影响？想一想长江、黄河为什么都是自西向东流？
\end{thinkbox}

\newpage

\section{\textcolor{orange!80!black}{第四课\quad 气候与天气}}

\subsection*{4.1\quad 气温与降水}

气温受\textbf{纬度}、\textbf{海陆位置}和\textbf{地形}三个因素影响：
\begin{itemize}
  \item 纬度越高气温越低（赤道热、两极冷）
  \item 沿海地区温差小，内陆温差大
  \item 海拔每升高 1000 米，气温下降约 $6^\circ$C
\end{itemize}

\subsection*{4.2\quad 世界气候类型}

\begin{center}
\begin{tabular}{|c|c|c|}
\hline
\textbf{气候带} & \textbf{特点} & \textbf{分布} \\ \hline
热带雨林气候 & 全年高温多雨 & 赤道附近（亚马孙、刚果盆地） \\
热带沙漠气候 & 全年高温少雨 & 南北回归线附近（撒哈拉沙漠） \\
温带季风气候 & 夏季高温多雨、冬季寒冷干燥 & 中国东部 \\
地中海气候 & 夏季炎热干燥、冬季温和多雨 & 地中海沿岸 \\
寒带气候 & 全年寒冷 & 极地地区 \\
\hline
\end{tabular}
\end{center}

\subsection*{4.3\quad 中国气候}

中国主要有\textbf{温带季风气候}、\textbf{亚热带季风气候}和\textbf{高原山地气候}。秦岭--淮河一线是重要的气候分界线：
\begin{itemize}
  \item 以北：温带季风气候，最冷月均温 $<0^\circ$C
  \item 以南：亚热带季风气候，最冷月均温 $>0^\circ$C
\end{itemize}

\begin{thinkbox}
"秦岭--淮河"线还是中国南方和北方的分界线。你能说出这条线以南和以北在饮食、住房、交通方面有什么不同吗？
\end{thinkbox}

\newpage

\section{\textcolor{orange!80!black}{第五课\quad 世界与中国地理概览}}

\subsection*{5.1\quad 大洲与大洋}

七大洲：亚洲、非洲、北美洲、南美洲、南极洲、欧洲、大洋洲 \\
四大洋：太平洋、大西洋、印度洋、北冰洋

\subsection*{5.2\quad 中国地理}

\begin{itemize}
  \item 面积：约 960 万平方公里（世界第三）
  \item 人口：约 14 亿
  \item 行政区划：23 个省、5 个自治区、4 个直辖市、2 个特别行政区
  \item 人口分布：东密西疏（胡焕庸线）
\end{itemize}

\begin{thinkbox}
胡焕庸线是中国人口地理的一条分界线——从黑河到腾冲。这条线以东集中了约 96\% 的人口。想一想：为什么这条线以西人口这么少？和地形、气候有什么关系？
\end{thinkbox}

\newpage

\section{\textcolor{orange!80!black}{第六课\quad 人口与城市}}

\subsection*{6.1\quad 世界人口分布}

世界人口主要集中在：
\begin{itemize}
  \item 东亚（中国东部、日本、韩国）
  \item 南亚（印度、孟加拉国）
  \item 欧洲西部
  \item 北美东部
\end{itemize}
这些地区的共同特点：地形平坦、气候适宜、靠近水源。

\subsection*{6.2\quad 城市化}

城市化是指人口从农村向城市集中的过程。中国的城市化率从 1978 年的约 18\% 上升到现在的约 65\%。

\begin{thinkbox}
城市和农村相比，各有什么优缺点？你能说出三个"农村好"的理由和三个"城市好"的理由吗？
\end{thinkbox}

\vspace{0.5cm}

\begin{center}
\textcolor{blue!70!black}{\large\textbf{—— 地理初步 · 完 ——}}
\end{center}

\end{document}
